\chapter*{Abstract}
\addcontentsline{toc}{chapter}{Abstract}

Chess players improve fastest when practice targets what they actually get wrong, yet most
puzzle platforms select puzzles by rating alone. This dissertation builds and evaluates the
\emph{Adaptive Chess Puzzle Advisor}, an end-to-end system that analyses a player's Chess.com games
with Stockfish, estimates tactical weaknesses across 23 motif categories, and serves puzzles from a
5.9-million-puzzle Lichess corpus and from positions mined from the player's own games. The
recommendation core is an online Bayesian item-response model, $P(\text{solve}) =
\sigma(\theta + \delta_k - b)$, which reduces exactly to the Glicko update in one dimension. It
uses Thompson Sampling on the category offsets $\delta_k$ to choose what to practise, and sets
puzzle difficulty so that the predicted solve probability is 0.65.

Tactic labelling, which every weakness estimate depends on, is done by \emph{PuzzleNet}: a multi-task
neural network trained here on 5.5 million Lichess puzzles, which reads a position and the line that
follows it and predicts the tactical category, the theme tags, and the difficulty together with an
uncertainty. Its rating likelihood treats each puzzle's own Lichess rating deviation as known
measurement noise, so the learned variance is the uncertainty the content leaves, which is the quantity
the Bayesian recommender needs for a puzzle nobody has played.

The evaluation is deliberately adversarial to the system's own assumptions. On a held-out sample, the
hand-written tagger agrees with Lichess theme labels on 58.4\% of puzzles (Cohen's $\kappa = 0.52$),
against 31.2\% ($\kappa = 0.23$) for the original single-move labeller; PuzzleNet, scored on exactly the
same puzzles, reaches 92.6\% ($\kappa = 0.91$), and $\kappa = 0.78$ in the harder condition the
application actually runs in, labelling engine lines rather than puzzle solutions. In a pre-registered replication, the original
bandit clearly beats random selection ($d = 5.4$) but fails two of its three pre-registered
criteria. Across 200 paired simulation runs the new policy cuts puzzles the player has less than a
30\% chance of solving from 44\% to 3\%. Whether this helps learning depends on the assumed learning
mechanism. On 300 real players stratified by rating, a player's rating band is predicted from how
they play with 59.1\% accuracy (five classes, chance 20\%, permutation $p = 0.002$). However, the
tactical categories a player will miss in future games are only marginally more predictable than
population base rates, and per-category weakness shows near-zero split-half reliability --- a result
that survives relabelling every one of those players' 40{,}067 critical positions with the neural
labeller, so it is not an artefact of crude labels.

The central conclusion is that the robust contribution is \emph{adaptive difficulty}, and that
personal tactical weaknesses are too subtle to diagnose from typical game histories and must be
learned online from targeted puzzle attempts. The generative-RL puzzle creation and human
evaluation envisaged in the project plan were not carried out. Both are discussed, with a
pre-registered study protocol for the latter.

\vspace{4mm}
\noindent\textbf{Keywords:} Thompson Sampling, item response theory, Glicko-2, adaptive learning,
chess tactics, player modelling, multi-task neural networks, heteroscedastic regression, simulation study.
